\chapter{Related Work}\label{chapter:related_work}

\section{Review Criteria}

Monitoring papers use the same terms for different goals, so each selected
approach is described using the following questions:

\begin{enumerate}
  \item What is its goal and what observation level does it use?
  \item Which ROS interfaces or other sources are observed?
  \item How are properties represented and checked?
  \item Is checking online, offline, or both?
  \item Where do collection and checking run?
  \item Which parts are configurable, extensible, or generated?
\end{enumerate}

The reviewed approaches observe a robot system at different depths of the
software stack. Throughout this chapter, three observation levels are
distinguished; Table~\ref{tab:observation-levels} defines each level and the
kind of information that is visible there.

\begin{table}[htbp]
\centering
\caption{The three observation levels used to classify related work.}
\label{tab:observation-levels}
\begin{tabular}{
  >{\raggedright\arraybackslash}p{.20\textwidth}
  >{\raggedright\arraybackslash}p{.33\textwidth}
  >{\raggedright\arraybackslash}p{.36\textwidth}}
\toprule
\textbf{Observation level} & \textbf{What is observed} &
\textbf{Visible information}\\
\midrule
Application level & The typed data that nodes exchange through ROS~2
interfaces. & Topic messages, service requests and responses, action feedback
and status, with ROS names and types preserved.\\
\addlinespace
Middleware and network level & The DDS middleware and the network traffic
below the ROS~2 API. & Discovery and endpoint events, DDS operations, packets,
and traffic flows; ROS meaning must be reconstructed.\\
\addlinespace
Execution and system level & Program execution inside the ROS~2 stack and the
operating system. & Callback execution, scheduling, tracepoints, and
operating-system events; timing rather than message content.\\
\bottomrule
\end{tabular}
\end{table}

\section{ROSMonitoring and ROSMonitoring 2.0}

ROSMonitoring is a topic-oriented runtime-verification framework for
ROS~1~\cite{ferrando2020rosmonitoring}. A YAML description selects the topics.
A generator creates monitor nodes and launch/remapping changes so selected
messages pass through a monitor. The monitor converts ROS messages to JSON
events. In online mode it sends events over WebSocket to an external oracle; in
offline mode it records a trace. An online filtering mode can block or
republish messages according to the oracle result. The property formalism is
therefore supplied by the external oracle rather than fixed by the ROS-facing
monitor.

ROSMonitoring 2.0 extends this approach to services and to explicitly ordered
topic observations~\cite{saadat2024rosmonitoring2}. Its Runtime Monitoring
Language supports event-pattern properties, including past-time Metric Temporal
Logic. The reported ROS~2 support is partial, and ROS~2 action monitoring is
not covered.

\section{ROMoSu}

ROMoSu is a configurable platform for runtime monitoring of ROS-based
applications~\cite{stadler2022towards,stadler2023romosu}. It separates a
monitoring-configuration phase from data collection. Reusable system-under-
monitoring types and monitor configurations describe topic fields, sampling,
validation, storage, and presentation. At runtime, a ROS connector and
adaptation manager collect data, and a runtime-data broker connects validation,
persistency, dashboards, and external constraint services.

The reported implementation focuses on ROS~1 topic data. Properties may be
checked by connected validation or complex-event-processing services. The
platform provides an administration interface and dashboard. Its main
variability lies in reusable configurations, selected data fields and rates,
and external services; code generation is not its central mechanism.

\section{Requirements-to-Monitor Generation for ROS 2}

Perez et al.\ present a requirements-to-monitor toolchain for
ROS~2~\cite{perez2022monitoringros2}. FRET, the Formal Requirements Elicitation
Tool, lets a user write structured FRETISH requirements. Ogma is a generator
that translates these requirements and the selected ROS~2 inputs into monitor
artifacts. Copilot in this work is a stream-based runtime-verification language
and framework, not an AI programming assistant. Ogma generates Copilot code,
which is compiled to C99 and packaged as a ROS~2 node.

The approach observes selected topic values at the application level and checks
generated temporal properties online. Its principal automation is the
translation from structured requirements to executable monitor code and ROS~2
packaging. Deployment variation beyond the generated package is not the main
focus reported in the paper.

\section{RTAMT}

RTAMT is a library for online and offline robustness monitoring of Signal
Temporal Logic (STL) and related specifications~\cite{yamaguchi2024rtamt}. A
property is evaluated over timestamped signals, producing a Boolean or
quantitative robustness result; the quantitative result shows how strongly a
signal satisfies or violates the property. The library supports discrete-time
and dense-time interpretations and provides Python and C++ evaluation
back-ends. Adapters can supply signals from cyber-physical and robotic
applications, including ROS.

RTAMT is primarily a checking engine. Configuration concerns the property, its
parser, and the evaluation interface. Source collection, routing, storage, and
multi-host deployment are supplied by the surrounding application rather than
by a complete monitoring platform, and the library does not prescribe where
its host process runs.

\section{ros2\_tracing}

\texttt{ros2\_tracing} is a low-overhead tracing framework for
ROS~2~\cite{bedard2022ros2tracing}. LTTng (Linux Trace Toolkit: next
generation) records tracepoints, which are predefined event probes in ROS~2,
the middleware layer, and user code. The resulting traces reveal callback
execution, message flow, latency, and scheduling.

This is execution- and middleware-level observation rather than a typed
application-property pipeline. It supports live trace collection and offline
analysis. Tracepoint selection and analysis are configurable; the reported
framework does not prescribe one property formalism or verdict service.

\section{Network and Security Monitoring}

ROS-FM performs fast monitoring at the network level using extended Berkeley
Packet Filter (eBPF) and eXpress Data Path (XDP) mechanisms
~\cite{rivera2020rosfm}. It observes traffic without adding a normal ROS
subscriber for every topic, applies security-oriented processing, and supports
distributed visualization. The trade-off is that typed application meaning has
to be recovered from lower-level traffic.

ROSIDS23 and ROSPaCe are datasets for intrusion-detection research rather than
general monitoring infrastructures~\cite{degirmenci2023rosids23,puccetti2024rospace}.
They combine network and, in the case of ROSPaCe, ROS~2 and operating-system
evidence. They show that security properties may require several observation
levels, but do not define a reusable property-checking deployment workflow.

\section{Containerized and Distributed Verification}

Aldegheri et al.\ generate ROS-compliant monitors and package them in Docker
containers~\cite{aldegheri2021containerized}. The approach studies monitor
placement between edge and cloud resources, so the same verification component
can be moved without rebuilding it. Checking runs online against the observed
ROS communication. Configuration in the reported work concerns container
placement and the monitored interfaces; the property logic is fixed inside the
generated monitors. Its main contribution is portable verification deployment
rather than broad ROS interface collection.

Distributed runtime verification more generally studies systems in which
observations and checking are spread across components and
hosts~\cite{francalanza2018distributed}. It includes property decomposition,
partial observations, clock and ordering assumptions, communication cost, and
failures. Separating a collector from a central verifier is one deployment
form, but it is not equivalent to decentralising the property itself.

\section{Digital-Twin Runtime Verification}

Betzer et al.\ place runtime verification within a cloud-located digital twin
for an autonomous mobile robot under uncertainty
~\cite{betzer2024digitaltwinrv}. Robot state is transmitted through MQTT. The
digital twin maintains and extends this state, and monitors generated from
TeSSLa (Temporal Stream-based Specification Language) properties perform
runtime validation. The design considers monitor placement inside the twin or
in separate services and includes a path by which checking can influence
actuation. Collection, digital-twin state, TeSSLa checking, and feedback are
therefore distinct architectural responsibilities.

\section{Field-Based Testing}

Caldas et al.\ combine runtime verification with field-based testing for
ROS-based robots~\cite{caldas2024runtime}. Their guidance covers data
collection, isolation, monitoring reactions, and use of operational evidence.
It emphasizes that simulation cannot reproduce every field condition and that
monitoring artifacts must support a wider testing process. The work is process
guidance rather than one fixed observation and checking pipeline.

\section{A Domain Model of Software Monitoring}

Rabiser et al.\ analyse the software-monitoring domain beyond
robotics~\cite{rabiser2019domain}. Their study of resource and requirements
monitoring approaches produces a domain model and a reference architecture
with three areas: monitoring setup, monitoring execution, and monitoring
support. Recurring responsibilities include monitoring definition,
instrumentation and probes, information collection, filtering and
aggregation, routing, checking, and persistency. The authors validate the
reference architecture by showing that existing monitoring solutions can be
expressed as instantiations of it.

The work is domain-independent rather than ROS-specific: it does not fix an
observation mechanism, property formalism, or deployment. Its variability lies
in the domain model itself, whose variation points describe where concrete
monitoring solutions differ. This thesis uses the identified tasks to organize
Chapter~\ref{chapter:background} and, in Chapter~\ref{chapter:evaluation},
checks whether the proposed ROS~2 reference architecture covers the same
responsibilities.

\section{Focused Summary}

Table~\ref{tab:related-work-focused} summarizes the reviewed approaches with
the same criteria. A dash means that the publication does not make that
capability a principal part of the approach; it does not imply that the
capability is impossible.

\begin{small}
\setlength{\tabcolsep}{3pt}
\begin{longtable}{
  >{\raggedright\arraybackslash}p{.16\textwidth}
  >{\raggedright\arraybackslash}p{.18\textwidth}
  >{\raggedright\arraybackslash}p{.21\textwidth}
  >{\raggedright\arraybackslash}p{.18\textwidth}
  >{\raggedright\arraybackslash}p{.18\textwidth}}
\caption{Focused comparison of selected monitoring approaches.}
\label{tab:related-work-focused}\\
\toprule
\textbf{Approach} & \textbf{Observation} & \textbf{Property and checking} &
\textbf{Execution and placement} & \textbf{Configuration and automation}\\
\midrule
\endfirsthead
\toprule
\textbf{Approach} & \textbf{Observation} & \textbf{Property and checking} &
\textbf{Execution and placement} & \textbf{Configuration and automation}\\
\midrule
\endhead
ROSMonitoring / 2.0 & Topics; 2.0 adds services and ordering &
External oracle; 2.0 also provides RML/past-time MTL &
Online or offline; inserted monitor nodes &
YAML selects sources; generator creates monitor/launch artifacts\\
\midrule
ROMoSu & Selected ROS~1 topic fields &
Validation and external services, including CEP &
Broker-connected platform services &
Reusable system and monitor configurations; UI\\
\midrule
FRET--Ogma--Copilot & Selected ROS~2 topic values &
Generated stream/temporal monitors &
Compiled monitor in a generated ROS~2 package &
Generation from structured requirements\\
\midrule
RTAMT & Timestamped signals supplied by an adapter &
STL robustness and Boolean checking &
Online or offline library &
Property and parser/API configuration\\
\midrule
\texttt{ros2\_tracing} & ROS~2, middleware, and system tracepoints &
Trace analysis; no fixed property language &
Low-overhead collection and offline analysis &
Tracepoint and analysis configuration\\
\midrule
ROS-FM & Network traffic &
Security rules and traffic analysis &
Distributed probe and visualization setup &
Probe and rule configuration\\
\midrule
Digital-twin RV & MQTT robot state and digital-twin state &
Generated TeSSLa monitors &
Cloud twin/service with optional feedback &
External properties and monitor-placement choices\\
\midrule
Rabiser et al. & Domain-independent probes and resources &
Requirements and resource checking &
Reference architecture with setup, execution, and support &
Domain model and variation points\\
\bottomrule
\end{longtable}
\end{small}

\section{Research Gaps}\label{sec:research-gap}

The review shows complementary strengths rather than one missing checking
language. The gaps relevant to this thesis are:

\begin{itemize}
  \item Typed ROS~2 topic observation is common, but the reviewed approaches do
  not jointly cover topic messages, service request/response events, and action
  feedback/status through one representation.

  \item Observation, processing, routing, checking, and deployment are often
  coupled to one platform or toolchain. Replacing a checker or changing its
  placement can therefore require new integration code.

  \item Configurability is usually concentrated in one area, such as selecting
  topics, writing properties, or placing generated monitors. Few approaches
  expose source, processing, checker, transport, output, and placement choices
  through one coherent model.

  \item Automation is available for particular transformations, especially
  requirements-to-monitor generation, but routine observation and multi-runtime
  wiring still recur across solutions.

  \item Evidence needs a stable identity and context if an online verdict is to
  be traced back to heterogeneous observations or replayed later.
\end{itemize}

These gaps motivate a reference architecture for the infrastructure between
ROS~2 observations and property-specific checkers. The proposed design is
presented only after this neutral review, in the next chapter.
